% ------------------------------------------------------------------------------
% 5 Results II -- Subgroups, redesign, re-enacted monitoring
% ------------------------------------------------------------------------------
\section{Results II: Subgroups, Trial Redesign, and Re-enacted
Monitoring}\label{sec:res-design}

\subsection{Subgroup analyses}

Figure~\ref{fig:subgroup} presents the seven pre-specified subgroup
analyses for the primary contrast. The hazard ratio is remarkably stable
across subgroups---every point estimate lies between $0.46$ and $0.81$
and every confidence interval overlaps the overall estimate of $0.62$.
Only one nominal interaction reaches $p<0.05$: sex
(interaction $p = 0.043$), with a substantially larger effect in men
(HR $0.46$, 95\% CI $0.33$--$0.64$) than in women (HR $0.79$, 95\% CI
$0.58$--$1.07$). Against the Bonferroni-adjusted threshold of
$0.0071$ appropriate for seven tests, this signal does not survive, and
no biological mechanism has been established for sex-specific efficacy
of 5-FU-based regimens in this setting. The honest reading is
\emph{consistent relative effect across all examined subgroups}, a
textbook illustration of why subgroup forest plots without
multiplicity-aware interpretation manufacture false discoveries.

\begin{figure}[htbp]
  \centering
  \includegraphics[width=0.82\textwidth, max height=0.44\textheight]{fig_subgroup_forest.png}
  \caption{Pre-specified subgroup analyses for disease-free survival,
  levamisole plus 5-FU versus observation. Right column: hazard ratios
  with 95\% CIs and interaction-test $p$-values (adjusted for nodal
  count).}
  \label{fig:subgroup}
\end{figure}

\subsection{The redesigned trial}

Table~\ref{tab:design} summarizes the fixed design and the
group-sequential redesign. For the target hazard ratio of $0.65$ at
one-sided $\alpha = 0.025$ with 90\% power, the fixed design requires
226 DFS events (approximately 322 patients under the assumed accrual
and follow-up). The three-look sequential design inflates the maximum
to 235 events but compensates with expected savings whenever the effect
is real or absent. The O'Brien--Fleming efficacy boundary is naturally
conservative early ($Z = 2.963$, one-sided $p = 0.0015$ at the first
interim) and relaxes toward $2.014$ at the final analysis; the
non-binding futility boundary starts at $Z = 0.079$, meaning a first
look with essentially no observed trend would trigger the futility
signal.

\begin{table}[htbp]
\centering\footnotesize
\caption{The redesigned group-sequential trial: boundaries and operating
characteristics from \texttt{gsDesign}. Efficacy uses O'Brien--Fleming
alpha-spending; futility uses a non-binding Hwang--Shih--DeCani
beta-spending ($\gamma = -4$).}
\label{tab:design}
\begin{tabular*}{\textwidth}{@{\extracolsep{\fill}} l c c c c c c}
\toprule
Analysis & Events & Effic.\ $Z$ & 1-sided $p$ & HR at bound & Futil.\ $Z$ & HR at bound \\
\midrule
Interim 1 (50\%)  & 118 & 2.963 & 0.0015 & 0.579 & 0.078 & 0.986 \\
Interim 2 (75\%)  & 177 & 2.359 & 0.0091 & 0.701 & 1.006 & 0.859 \\
Final (100\%)     & 235 & 2.014 & 0.0220 & 0.769 & 2.014 & 0.769 \\
\midrule
\multicolumn{7}{@{}l@{}}{\emph{Design-theoretic probabilities:}
$P(\text{reject} \mid \mathrm{HR}=0.65) = 0.900$;\;
$P(\text{reject} \mid \mathrm{HR}=1) = 0.024$ (futility adhered)} \\
\bottomrule
\end{tabular*}
\end{table}

Figure~\ref{fig:gsbound} visualizes the boundaries on the $Z$ and
one-sided $p$-value scales, and Figure~\ref{fig:power} shows how power
and required events vary with the target hazard ratio: the design
maintains 90\% power at $\mathrm{HR}=0.65$ with 226 events, but power
falls steeply for weaker effects---at $\mathrm{HR}=0.80$, nearly 850
events would be required, a useful reminder of why ambitious trials
with modest expected effects grow so large.

\begin{figure}[htbp]
  \centering
  \includegraphics[width=0.80\textwidth, max height=0.38\textheight]{fig_gs_boundaries.png}
  \caption{Group-sequential boundaries of the redesigned trial on the
  $Z$ scale (left axis) and one-sided $p$-value scale (right axis).
  Red: efficacy (O'Brien--Fleming spending); blue: non-binding
  futility (Hwang--Shih--DeCani).}
  \label{fig:gsbound}
\end{figure}

\begin{figure}[htbp]
  \centering
  \includegraphics[width=0.78\textwidth, max height=0.36\textheight]{fig_power_curve.png}
  \caption{Power as a function of the target hazard ratio under the
  Schoenfeld--Freedman approximation (one-sided $\alpha=0.025$, 1:1
  allocation), annotated with the required event count at
  $\mathrm{HR}=0.65$.}
  \label{fig:power}
\end{figure}

\subsection{Re-enacting interim monitoring on the real data}

The most striking result of this project emerges when the redesigned
monitoring plan is applied to the \emph{actual} event stream.
Censoring every patient's follow-up at the time at which the 118th DFS
event was observed (information fraction 50\%) and refitting the two-arm Cox model
yields an interim hazard ratio of $0.52$ and a $Z$-statistic of
$3.38$---\emph{above the efficacy boundary of $2.963$}
(Figure~\ref{fig:reenact}). Under the redesign, the trial would have
stopped at the first interim analysis and declared the levamisole plus
5-FU regimen superior to observation, after only 118 of the 506 events
that were eventually observed.

Conditional power calculations quantify how decisive this interim
result is: the probability of eventually crossing the final boundary is
70.6\% even under the null, 99.8\% under the design alternative, and
essentially 100\% if the observed interim trend persists. A cross-check
with \texttt{gsDesign}'s own conditional-power machinery
(\texttt{gsCP}) returns somewhat higher values ($0.834$ and $0.999$ under
the null and alternative), because it accumulates the crossing
probabilities at \emph{both} remaining analyses with boundaries
re-optimized after the interim, whereas the hand-rolled calculation uses
a single-look normal approximation to the final analysis only.
The contrast with the Bayesian view is instructive: sequential
posterior probabilities under the skeptical prior
(Section~\ref{sec:res-bayes}) would have signalled benefit even
earlier, at the 25\% information look.

\begin{figure}[htbp]
  \centering
  \includegraphics[width=0.80\textwidth, max height=0.38\textheight]{fig_gs_reenactment.png}
  \caption{Re-enacted interim monitoring using the real trial data.
  The observed interim result after 118 events (HR $0.52$,
  $Z = 3.38$, green marker) crosses the O'Brien--Fleming efficacy
  boundary ($2.96$), so the redesigned trial would have stopped for
  efficacy at the first interim analysis.}
  \label{fig:reenact}
\end{figure}

Two interpretive cautions apply. First, the historical trial was not
designed with this monitoring plan, so the re-enactment is a
counterfactual exercise, not a critique of the original conduct;
its message is about how much evidence was already present early in
the event stream. Second, the interim hazard ratio of $0.52$ is more
extreme than the final estimate of $0.62$, a reminder that early
random highs in $Z$ are partly regression artefacts; the design's
conservative early boundary exists precisely to price this randomness
in.
